% MAPA artigo_3:5.4
\begin{proposition}[forma ancorada e comprimento estrutural]\label{map:artigo_3:5.4}
\label{artigo_3:prop:anchored-relu-code}
Se $g\in\Aclass^\square_{N,H}([0,1])$, sua forma reduzida possui
$M\le N-1$ hinges e pode ser escrita, tomando o centro $x_*$ da folha mais à
esquerda, como
\begin{equation}
 g(x)=g(x_*)+\sigma_0(x-x_*)
      +\sum_{j=1}^{M}(\sigma_j-\sigma_{j-1})\rho(x-\tau_j).
 \label{artigo_3:eq:anchored-relu-form}
\end{equation}
Os valores $g(x_*)$ e $\sigma_i$ têm altura no máximo $H$, os saltos têm
altura no máximo $2H^2$, e os pontos diádicos são transmitidos por seus pares
estruturais. Consequentemente, existe uma codificação prefix-free da forma
ancorada com
\begin{equation}
 L_{\rm ReLU}^{\rm anc}(g)=O(N\log(2H)).
 \label{artigo_3:eq:anchored-network-bits}
\end{equation}

Na direção contrária, se os coeficientes de
\eqref{artigo_3:eq:reduced-relu-reverse} têm altura no máximo $W$ e as quebras têm
profundidade no máximo $D$, então
\begin{equation}
 \log_2 H_{\rm can}(R)
 \le C\bigl[M\log(2W)+D+\log(M+1)\bigr]
 \label{artigo_3:eq:reverse-height-bit-bound}
\end{equation}
para uma constante universal $C$.
\end{proposition}

\begin{proof}
Na ida, cada salto é a diferença de duas inclinações de altura no máximo
$H$, logo sua altura não excede $2H^2$. A altura atômica limita os inteiros
estruturais da célula, e todos os campos usam $O(\log(2H))$ bits. Somar os
$O(N)$ campos prova \eqref{artigo_3:eq:anchored-network-bits}. O ancoramento evita
expandir desnecessariamente $g(x_*)-\sigma_0x_*$ em um intercepto com
denominador diádico explícito.

Na volta, cada inclinação local é soma de no máximo $M+1$ racionais. Um
denominador comum tem no máximo $O(M\log(2W))$ bits. Avaliar em um centro de
profundidade no máximo $D$ acrescenta $D+O(\log(M+1))$ bits. Exponenciar a
cota de comprimento fornece \eqref{artigo_3:eq:reverse-height-bit-bound}.
\end{proof}
